\section{Residual transports and products along different paths}
\label{app:residual-coherence}

Section~\ref{sec:gkt-decoration} compares the diagram $\mathfrak D$ with GKT
wall-crossing one ray at a time: on each shifted ray, the factor carried by
the wall of $\mathfrak D$ is the product of the GKT factor and the residual
factor (Definition~\ref{def:raywise-logarithms}).  This appendix extends the
comparison to products of transported functions.  A function in the chart
algebra of one region of $\mathfrak D$ can be carried to the chart algebra of
another region in two ways: by the transport through $\mathfrak D$, which
applies the slab transitions and the wall factors met along a path, and by
the GKT transport, which applies the same slab transitions and, in place of
each wall factor, the GKT factor on its ray.  The two results differ by one
automorphism of the target chart algebra, the residual transport of the path
(Definition~\ref{def:gkt-transport}); being an algebra automorphism, it also
relates the products of functions carried along the same path.  For functions carried
along different paths the residual transports differ in general, even when
the paths have the same endpoints.  We therefore correct each factor by an
automorphism of the target algebra, so that the corrected factors differ from
their GKT transports by one common automorphism: for two factors the
correction uses the midpoint of the two residual transports
(Lemma~\ref{prop:residual-midpoint}), and for more factors an interpolation
weighted by degrees (Definition~\ref{def:corrected-product}).  The corrected
products for different parenthesizations are related by explicit
automorphisms that satisfy the cocycle identity and, for three factors, are
not the identity in general (Proposition~\ref{thm:residual-coherence}); a
change of the paths acts through the changes of the GKT transports of the
factors and of the interpolated residual transport
(Proposition~\ref{prop:corridor-groupoid}).

\subsection{Transports along a path}

Fix a labelled truncation (Section~\ref{subsec:raywise-factorization}), a
deformation order, and a region of $\mathfrak D$
(Definition~\ref{def:Jacobi-wall}), the \emph{target region}, and work in the
corresponding finite nilpotent quotient.  The paths below are transverse to
the walls, realize admissible paths
(Definition~\ref{def:framed-path}), and end in the target region.

\begin{definition}[Transports along a path]
\label{def:gkt-transport}
Let $\gamma$ be such a path.
\begin{enumerate}
\item The \emph{transport through $\mathfrak D$} along $\gamma$ is the
composite $P_\gamma$, in the order in which they are met, of the slab
transitions and of the contributions $\Theta_{\mathfrak d}^{\pm1}$ of the
walls crossed by $\gamma$ (Definition~\ref{def:Jacobi-wall}).  It carries the
chart algebra at the initial point of $\gamma$ to the chart algebra of the
target region.  In the notation of Section~\ref{subsec:KS-classes},
$P_\gamma=P_{\mathfrak p}$ for the admissible path $\mathfrak p$ that $\gamma$
realizes.
\item The \emph{GKT transport} $P^{\GKT}_\gamma$ is obtained from $P_\gamma$
by replacing, for every crossed wall $\mathfrak d$ whose factor lies in the
ray group of a shifted ray $\mathfrak r$, the contribution
$\Theta_{\mathfrak d}$ of a positive crossing by $\exp(G_{\mathfrak r})$ and
the contribution $\Theta_{\mathfrak d}^{-1}$ of a negative crossing by
$\exp(-G_{\mathfrak r})$, where $G_{\mathfrak r}$ is the GKT logarithm of
Definition~\ref{def:raywise-logarithms}; at a chord, $\mathfrak r$ is the
chord ray, on either half of the chord.  The slab transitions are not
changed.
\item The \emph{residual transport} of $\gamma$ is
\[
 \mathcal Q_\gamma=P_\gamma\bigl(P^{\GKT}_\gamma\bigr)^{-1},
\]
an automorphism of the chart algebra of the target region, congruent to the
identity modulo the ideal generated by the $t_i$.
\end{enumerate}
\end{definition}

The GKT transport depends on the path and not only on its endpoints, and it
is not in general a GKT transition.  For the two labels of twists $(1,0)$
and $(0,1)$ of Proposition~\ref{prop:first-decoration}, the paths
$\gamma_+$ and $\gamma_-$ of Proposition~\ref{prop:two-label-transport} run
in one chart from $\ell$ to $r$, and their transports through $\mathfrak D$ are the path-ordered products
$\Theta_{\gamma_+}=\Theta_{\gamma_-}=\Sigma_1\Sigma_2$.  The GKT transport
along $\gamma_+$ is the transition \eqref{eq:mixed-GKT-transition},
$\Sigma_2Z_{1/4}\Sigma_1=Z_1\Sigma_1\Sigma_2$, and that along $\gamma_-$ is
$\Sigma_1\Sigma_2$.  With $Z_c$ as in \eqref{eq:mixed-total-factor}, the
residual transports of $\gamma_+$ and $\gamma_-$ are therefore $Z_{-1}$ and
the identity, although the two paths have the same endpoints.  If the system
contains two labels $j,j'$ of twist $\alpha$, the GKT transport along $\gamma_+$ omits the GKT factor
$\exp(\frac12t_jt_{j'}X_{2,0})$, whose ray carries no wall.  On the negative
half-chords the wall factor and $\exp(G_{\mathfrak r})$ agree for labelled
systems with at most three labels, but not in general: for a support $J$ of
type $(3\alpha,\beta)$ the GKT factor on the ray of $(2,1)$ has the
component $\frac{15}8t_JX_{3,1}$, whereas the negative half-chord factor has
no $t_J$-component and the positive half-chord factor has the component
$-\frac38t_JX_{3,1}$.

\subsection{Two factors}

Let $\gamma_1,\gamma_2$ be two paths as above, let $f_i$ be a function in
the chart algebra at the initial point of $\gamma_i$, let $\mu$ be the
multiplication of the chart algebra of the target region, and write
$\mathcal Q_i=\mathcal Q_{\gamma_i}$.  Replacing the wall factors by GKT
factors changes the transported function $P_{\gamma_i}f_i$ into
$P^{\GKT}_{\gamma_i}f_i=\mathcal Q_i^{-1}P_{\gamma_i}f_i$.  If
$\mathcal Q_1=\mathcal Q_2=\mathcal Q$, then
$\mu(P_{\gamma_1}f_1,P_{\gamma_2}f_2)
=\mathcal Q\,\mu(P^{\GKT}_{\gamma_1}f_1,P^{\GKT}_{\gamma_2}f_2)$, so the two
products differ by one automorphism of the target algebra.  In general
$\mathcal Q_1\neq\mathcal Q_2$, and we correct the two factors so that they
differ from their GKT transports by one common automorphism.  Rational powers
of automorphisms congruent to the identity are the finite polynomials
$g^c=\exp(c\log g)$.  The \emph{midpoint} of $\mathcal Q_1$ and
$\mathcal Q_2$ and the \emph{corrections} of the two factors are
\begin{equation}\label{eq:residual-midpoint}
 H=\mathcal Q_1(\mathcal Q_1^{-1}\mathcal Q_2)^{1/2},
 \qquad C_i=H\mathcal Q_i^{-1}.
\end{equation}
Thus $H$ is the value at $c=\frac12$ of the one-parameter family
$c\mapsto\mathcal Q_1(\mathcal Q_1^{-1}\mathcal Q_2)^c$, which runs from
$\mathcal Q_1$ at $c=0$ to $\mathcal Q_2$ at $c=1$.

\begin{lemma}[Midpoint correction]\label{prop:residual-midpoint}
The midpoint is symmetric in $\mathcal Q_1,\mathcal Q_2$, and
\begin{equation}\label{eq:corrected-product}
 C_iP_{\gamma_i}=HP^{\GKT}_{\gamma_i},
 \qquad
 \mu(C_1P_{\gamma_1}f_1,C_2P_{\gamma_2}f_2)
 =H\mu(P^{\GKT}_{\gamma_1}f_1,P^{\GKT}_{\gamma_2}f_2).
\end{equation}
It is natural under finite restriction.  If both paths are continued by a
common path from the target region, with transport $\Psi$ through
$\mathfrak D$ and GKT transport $\Psi'$, then $\mathcal Q_i$ is replaced by
$\Psi\mathcal Q_i\Psi'^{-1}$ and $H$ by $\Psi H\Psi'^{-1}$.
\end{lemma}

\begin{proof}
With $\phi=\mathcal Q_1^{-1}\mathcal Q_2$,
$\mathcal Q_2(\mathcal Q_2^{-1}\mathcal Q_1)^{1/2}
=\mathcal Q_1\phi(\phi^{-1})^{1/2}=\mathcal Q_1\phi^{1/2}$, proving symmetry.
The first identity in \eqref{eq:corrected-product} follows from
$\mathcal Q_i=P_{\gamma_i}(P^{\GKT}_{\gamma_i})^{-1}$; the second uses that
$H$ is an algebra automorphism.  After continuation, $P_{\gamma_i}$ and
$P^{\GKT}_{\gamma_i}$ become $\Psi P_{\gamma_i}$ and
$\Psi'P^{\GKT}_{\gamma_i}$, so $\mathcal Q_i$ becomes
$\Psi\mathcal Q_i\Psi'^{-1}$ and $\mathcal Q_1^{-1}\mathcal Q_2$ becomes
$\Psi'\mathcal Q_1^{-1}\mathcal Q_2\Psi'^{-1}$.  Rational powers commute with
conjugation, since $\log$ and $\exp$ do, which gives the formula for $H$.
\end{proof}

For the paths $\gamma_1=\gamma_+$ and $\gamma_2=\gamma_-$ above and functions
$f_1,f_2$ in the chart algebra at $\ell$, the midpoint is $H=Z_{-1/2}$, the corrected factors
are $Z_{1/2}\Sigma_1\Sigma_2f_1$ and $Z_{-1/2}\Sigma_1\Sigma_2f_2$, and their
product is the image under $Z_{-1/2}$ of the product
$(Z_1\Sigma_1\Sigma_2f_1)(\Sigma_1\Sigma_2f_2)$ of the GKT transports.
Without the corrections, the product $(\Sigma_1\Sigma_2f_1)(\Sigma_1\Sigma_2f_2)$
of the transports through $\mathfrak D$ is not the image of that product
under any automorphism of the target algebra independent of $f_1$ and $f_2$,
by the last assertion of Proposition~\ref{prop:corridor-groupoid}.

\subsection{Several factors}

Let $\gamma_1,\dots,\gamma_n$ be paths as above, with functions $f_i$ and
residual transports $\mathcal Q_i=\mathcal Q_{\gamma_i}$, and let each input
$i$ carry a positive degree $r_i$.  For positive degrees $r,s$ and
automorphisms $g,h$ congruent to the identity, set
\begin{equation}\label{eq:weighted-interpolation}
 M_{r,s}(g,h)=g(g^{-1}h)^{s/(r+s)},
\end{equation}
the value at $c=s/(r+s)$ of the family $c\mapsto g(g^{-1}h)^c$ from $g$ to
$h$; thus
$M_{r,r}(\mathcal Q_1,\mathcal Q_2)=H$.  Then $M_{r,s}(g,h)=M_{s,r}(h,g)$,
and $M_{r,s}(\Psi g\Psi'^{-1},\Psi h\Psi'^{-1})=\Psi M_{r,s}(g,h)\Psi'^{-1}$
for all $\Psi,\Psi'$ for which both sides are defined, by the computation in
the proof of Lemma~\ref{prop:residual-midpoint}.

\begin{definition}[Interpolants and corrected products]
\label{def:corrected-product}
A \emph{product tree} $\mathsf T$ on the inputs $1,\dots,n$ is a planar
rooted binary tree whose leaves are $1,\dots,n$ in this order, that is, a
parenthesization of a product of $n$ factors; each vertex carries the sum of
the degrees of the leaves below it.
\begin{enumerate}
\item The \emph{interpolant} $H_{\mathsf T}$ is defined recursively:
$H_{\mathsf T}=\mathcal Q_i$ if $\mathsf T$ is the leaf $i$, and
$H_{\mathsf T}=M_{r,s}(H_{\mathsf T_1},H_{\mathsf T_2})$ if the root of
$\mathsf T$ has subtrees $\mathsf T_1,\mathsf T_2$ of degrees $r,s$.
\item The \emph{corrected product} $\mathsf P_{\mathsf T}$ is defined by the
same recursion: $\mathsf P_{\mathsf T}=P_{\gamma_i}f_i$ if $\mathsf T$ is the
leaf $i$, and
\[
 \mathsf P_{\mathsf T}
 =\mu\bigl(H_{\mathsf T}H_{\mathsf T_1}^{-1}\mathsf P_{\mathsf T_1},\,
 H_{\mathsf T}H_{\mathsf T_2}^{-1}\mathsf P_{\mathsf T_2}\bigr).
\]
\end{enumerate}
\end{definition}

At each vertex, the two products of the subtrees are corrected by the
automorphisms $H_{\mathsf T}H_{\mathsf T_j}^{-1}$ before they are
multiplied.  For two leaves of equal degree, $H_{\mathsf T}=H$ and these
corrections are the $C_i$ of \eqref{eq:residual-midpoint}.  The degrees make
the interpolant independent of the tree when the residual transports
commute: if the $\mathcal Q_i$ commute, induction on $\mathsf T$ gives
\[
 \log H_{\mathsf T}=\frac{\sum_ir_i\log\mathcal Q_i}{\sum_ir_i}
\]
for every product tree $\mathsf T$.

\begin{proposition}[Degree-weighted coherence]
\label{thm:residual-coherence}
For every product tree $\mathsf T$,
\[
 \mathsf P_{\mathsf T}=H_{\mathsf T}\prod_iP^{\GKT}_{\gamma_i}f_i.
\]
For product trees $\mathsf T,\mathsf T',\mathsf T''$ on the same inputs, the
\emph{tree transition}
$\mathcal A_{\mathsf T\leftarrow\mathsf T'}=H_{\mathsf T}H_{\mathsf T'}^{-1}$
carries $\mathsf P_{\mathsf T'}$ to $\mathsf P_{\mathsf T}$, and
\[
 \mathcal A_{\mathsf T\leftarrow\mathsf T'}
 \mathcal A_{\mathsf T'\leftarrow\mathsf T''}
 =\mathcal A_{\mathsf T\leftarrow\mathsf T''}.
\]
In particular the composite of the tree transitions around any cycle of
trees is the identity; for four inputs this includes the pentagon of the
associahedron.  Let $\mathcal U$ be the quotient of the free associative
algebra $\QQ\langle x,y,z\rangle$ by the span of the words containing some
letter twice, let $\mathcal Q_1=e^x$, $\mathcal Q_2=e^y$, $\mathcal Q_3=e^z$
in $\mathcal U$ have equal degrees, and let $\mathsf T_\ell=((1,2),3)$ and
$\mathsf T_r=(1,(2,3))$.  Then
\begin{equation}\label{eq:first-associator}
 \log\mathcal A_{\mathsf T_\ell\leftarrow\mathsf T_r}
 =\log\bigl(H_{\mathsf T_\ell}H_{\mathsf T_r}^{-1}\bigr)
 =\frac1{72}[[x,z],y].
\end{equation}
In particular $\mathcal A_{\mathsf T_\ell\leftarrow\mathsf T_r}$ is not the
identity, and the two corrected triple products differ in general.
\end{proposition}

\begin{proof}
At an internal vertex, the corrections carry both products of the subtrees
to the same interpolant $H_{\mathsf T}$.  Since this interpolant is an
algebra automorphism, induction gives the displayed form of
$\mathsf P_{\mathsf T}$.  Ratios of interpolants carry one tree product to
another, and cancellation proves the cocycle identity.
Baker--Campbell--Hausdorff expansion for three equal degrees gives the
common linear term $(x+y+z)/3$, no quadratic difference, and, modulo words
containing a letter twice, the cubic difference
\[
 \frac1{72}(xzy-zxy-yxz+yzx)
 =\frac1{72}[[x,z],y],
\]
which is \eqref{eq:first-associator}.  By comparison, unweighted iteration
of the midpoint assigns leaf weights $(1/4,1/4,1/2)$ and $(1/2,1/4,1/4)$ in
the two parenthesizations already in the abelian quotient.
\end{proof}

Without the square-zero condition,
$\log\mathcal A_{\mathsf T_\ell\leftarrow\mathsf T_r}$ has further cubic
terms, such as multiples of $[x,[x,y]]$.

\subsection{Change of paths}

A change of the paths acts on the individual factors, which we call the
\emph{legs} of the product, as well as on the interpolant.  Let
$\boldsymbol\gamma=(\gamma_i)$ and $\boldsymbol\gamma'=(\gamma'_i)$ be two
tuples of paths as above such that $\gamma_i$ and $\gamma'_i$ have the same
endpoints for every $i$.  For a fixed product tree, write
$H_{\boldsymbol\gamma}$ and $H_{\boldsymbol\gamma'}$ for the interpolants of
the two tuples, and $H_{\boldsymbol\gamma,\mathsf T}$ when the tree
$\mathsf T$ is to be specified.

\begin{definition}[Leg maps and output map]
\label{def:leg-output-maps}
The \emph{leg maps} and the \emph{output map} of the change from
$\boldsymbol\gamma$ to $\boldsymbol\gamma'$ are the automorphisms of the
chart algebra of the target region
\begin{equation}\label{eq:corridor-maps}
 W_i^{\boldsymbol\gamma'\leftarrow\boldsymbol\gamma}
   =P^{\GKT}_{\gamma'_i}\bigl(P^{\GKT}_{\gamma_i}\bigr)^{-1},
 \qquad
 B^{\boldsymbol\gamma'\leftarrow\boldsymbol\gamma}
   =H_{\boldsymbol\gamma'}H_{\boldsymbol\gamma}^{-1}.
\end{equation}
The leg map $W_i^{\boldsymbol\gamma'\leftarrow\boldsymbol\gamma}$ changes
the GKT transport of the $i$th leg, and the output map changes the
interpolant.
\end{definition}

\begin{proposition}[Change of admissible paths]
\label{prop:corridor-groupoid}
For $a_i=P^{\GKT}_{\gamma_i}f_i$,
\begin{equation}\label{eq:corridor-square}
 H_{\boldsymbol\gamma'}\prod_iP^{\GKT}_{\gamma'_i}f_i
 =B^{\boldsymbol\gamma'\leftarrow\boldsymbol\gamma}H_{\boldsymbol\gamma}
   \prod_i\bigl(W_i^{\boldsymbol\gamma'\leftarrow\boldsymbol\gamma}a_i\bigr).
\end{equation}
The leg maps and output maps compose under successive changes of paths, and
every composite of output maps and of the tree transitions of
Proposition~\ref{thm:residual-coherence} from
$(\boldsymbol\gamma,\mathsf T)$ to $(\boldsymbol\gamma',\mathsf T')$ equals
$H_{\boldsymbol\gamma',\mathsf T'}H_{\boldsymbol\gamma,\mathsf T}^{-1}$.

The leg maps cannot in general be absorbed into one automorphism of the
target algebra: if one leg changes by a nontrivial unital automorphism $W$
and another leg is fixed, no automorphism $\psi$ of the target algebra
satisfies $\psi(ab)=W(a)b$ for all $a,b$.
\end{proposition}

\begin{proof}
The identities
$P^{\GKT}_{\gamma'_i}f_i=W_i^{\boldsymbol\gamma'\leftarrow\boldsymbol\gamma}a_i$
and $B^{\boldsymbol\gamma'\leftarrow\boldsymbol\gamma}H_{\boldsymbol\gamma}
=H_{\boldsymbol\gamma'}$ give \eqref{eq:corridor-square}.  For a third tuple
$\boldsymbol\gamma''$, cancellation gives
\[
 W_i^{\boldsymbol\gamma''\leftarrow\boldsymbol\gamma'}
 W_i^{\boldsymbol\gamma'\leftarrow\boldsymbol\gamma}
 =W_i^{\boldsymbol\gamma''\leftarrow\boldsymbol\gamma},
 \qquad
 B^{\boldsymbol\gamma''\leftarrow\boldsymbol\gamma'}
 B^{\boldsymbol\gamma'\leftarrow\boldsymbol\gamma}
 =B^{\boldsymbol\gamma''\leftarrow\boldsymbol\gamma}.
\]
Tree transitions are likewise ratios of the corresponding interpolants, so
every composite telescopes to
$H_{\boldsymbol\gamma',\mathsf T'}H_{\boldsymbol\gamma,\mathsf T}^{-1}$.
Finally, if $\psi(ab)=W(a)b$, then $b=1$ gives $\psi=W$, whereas $a=1$ gives
$\psi=\mathrm{id}$.  Hence such a $\psi$ exists only when the leg change is
trivial.
\end{proof}
